% Adaptación del preámbulo de Paper III v1.5 español; no plantilla standalone de Pandoc.
\documentclass[english,11pt,a4paper]{article}
\usepackage[margin=1in]{geometry}
\usepackage{fontspec}
\setmainfont{lmroman10-regular.otf}[BoldFont=lmroman10-bold.otf,ItalicFont=lmroman10-italic.otf,BoldItalicFont=lmroman10-bolditalic.otf]
\setsansfont{lmsans10-regular.otf}[BoldFont=lmsans10-bold.otf]
\setmonofont{lmmono10-regular.otf}[Scale=MatchLowercase]
\newlength{\LeanCharWidth}
\usepackage{amsmath,amssymb}
\usepackage[english]{babel}
\usepackage{microtype,parskip}
\usepackage{graphicx,float,caption}
\usepackage{longtable,booktabs,array,calc}
\usepackage{xcolor}
\usepackage{xurl}
\usepackage{hyperref,bookmark}
\usepackage{etoolbox,needspace}
\setcounter{secnumdepth}{-\maxdimen}
\setlength{\emergencystretch}{3em}
\providecommand{\tightlist}{\setlength{\itemsep}{0pt}\setlength{\parskip}{0pt}}
\newcounter{none}
\AtBeginEnvironment{longtable}{\Needspace{8\baselineskip}\small}
\pretocmd{\section}{\Needspace{5\baselineskip}}{}{}
\pretocmd{\subsection}{\Needspace{4\baselineskip}}{}{}
\pretocmd{\subsubsection}{\Needspace{4\baselineskip}}{}{}
\captionsetup{font=small,labelfont=bf}
\urlstyle{same}
\hypersetup{colorlinks=true,linkcolor=blue,urlcolor=blue,citecolor=blue,pdflang={en-US}}
\title{Formalization companion to exact chordal stability}
\author{Juan Pablo Traverso Gianini}
\date{}\begin{document}\raggedbottom\maketitle
Preprint 1.0, 18 October 2026. Formal application of Papers IV and VI.
This companion identifies the final source package, its public statement
and its trusted boundary. The note develops the stated stability
consequence of the source papers; it does not claim a new solution of
Erdős Problem 81.

\subsection{Public statement and
scope}\label{public-statement-and-scope}

The public endpoint is
\nolinkurl{ChordalStabilityNote.eventual_root_stability}. It states that
there is an order threshold N such that every graph G on Fin n, with n
at least N, satisfying chordality,
\texttt{cp\ G\ \textless{}=\ Mval\ n}, and
\texttt{2000\ *\ (Mval\ n\ -\ cp\ G)\ \textless{}=\ n}, satisfies
\nolinkurl{RootStability} for the actual degree-defined set
\texttt{\{v\ \textbar{}\ n\ \textless{}\ 2\ *\ G.degree\ v\}}.

The explicit upper bound \texttt{cp\ G\ \textless{}=\ Mval\ n} makes
natural subtraction agree with the nonnegative deficit in the
manuscript. No structural witness, local modification bound, exact
partition formula, or root is assumed in this endpoint. These are
obtained in its proof. The threshold is the maximum of the
uniform-rounding threshold and \texttt{2\ *\ 10\^{}13}; the numerical
term does not replace the unknown rounding threshold.

\nolinkurl{RootStability} is a record of conclusions. It includes:

{\def\LTcaptype{none} % do not increment counter
\begin{longtable}[]{@{}
  >{\raggedright\arraybackslash}p{(\linewidth - 2\tabcolsep) * \real{0.5000}}
  >{\raggedright\arraybackslash}p{(\linewidth - 2\tabcolsep) * \real{0.5000}}@{}}
\toprule\noalign{}
\begin{minipage}[b]{\linewidth}\raggedright
Field
\end{minipage} & \begin{minipage}[b]{\linewidth}\raggedright
Mathematical conclusion
\end{minipage} \\
\midrule\noalign{}
\endhead
\bottomrule\noalign{}
\endlastfoot
\nolinkurl{isClique} & The degree-defined root is a clique \\
\nolinkurl{exactFormula} & \texttt{cp\ =\ B(n,p)\ -\ A\ -\ 2m} over
integers \\
\nolinkurl{four} & \texttt{cp\ =\ cpLe\ 4} \\
\nolinkurl{deficitFormula} & \texttt{M-cp\ =\ (M-B)+A+2m} \\
\nolinkurl{editIdentity} & The edit distance to the complete split
reference equals \nolinkurl{A+m} \\
\nolinkurl{editBound} & That edit distance is at most \nolinkurl{M-cp} \\
\nolinkurl{square} & \texttt{(6p-2n-1)\^{}2\ \textless{}=\ 24(M-cp)+25} \\
\nolinkurl{canonical} & The root is exactly the set determined by the
degree test \\
\end{longtable}
}

The integer formula for B avoids truncated subtraction. The squared
inequality is expressed over rationals. These choices preserve the
manuscript's arithmetic meaning.

\subsection{Physical adapters and proof
dependencies}\label{physical-adapters-and-proof-dependencies}

\nolinkurl{PartitionAdapter} converts the actual IV exact partition and the
completion of its physical packing into the VI definition of clique
partition. It identifies the literal exterior edges, core edges and
missing spokes with VI's counting functions. A clique root has
\texttt{Cdef\ =\ 0}. These identifications are proved, not introduced as
interface axioms.

\nolinkurl{WindowArithmetic} and \nolinkurl{GraphWindow} establish the scalar
entry bounds, the comparison between maximum modification degree and
total edit distance, and recognition of the core by degree.
\nolinkurl{PaidRootAssembly} then consumes a physical partition and its
paid ledger to invoke the actual parametric theorem
\nolinkurl{PaperVI.cp_window_K4} with parameter four. This intermediate
theorem has a paid-root hypothesis, whereas the public endpoint does
not.

\nolinkurl{NearAssembly} applies this intermediate result to
\nolinkurl{W.regularized.root} from IV's \nolinkurl{NearStructureWitness}.
That distinction is essential: the physical accounts belong to the
regularized root, rather than to an unrelated field named core.
\nolinkurl{FarGap} rules out the far regime by comparing an actual rounded
packing's completed partition with the actual optimum \texttt{cp\ G}.
\nolinkurl{GlobalAssembly} supplies the certified fractional optimum,
invokes uniform rounding, obtains the near witness by localization, and
assembles the conclusion.

The selected proof chain requires IV's \nolinkurl{uniformTransferAt},
\nolinkurl{NearH1Localization.exists_nearStructureWitness_of_nearRegime},
its localization-by-descent dependency, the physical adapters, and VI's
parametric exact construction. It excludes
\nolinkurl{PaperIV.Erdos81Unconditional}. The application retains its
mathematical dependence on IV and VI while having its own source project
and endpoint. It does not use Paper V's transfer theorem or VI's
additional recourse results. The 18-module VI input cone is the frozen
parametric source cut used for this proof; the principal note cites the
current editorial manuscript without claiming that its additional
modules entered this parametric source cone.

\subsection{Final validation and source
package}\label{final-validation-and-source-package}

The project uses \nolinkurl{leanprover/lean4:v4.35.0} with Mathlib pinned
to \nolinkurl{56a005855ea17e0b0958f73ec801ae140b0a1679}. Its source closure
has 343 modules: eight modules of this note, 317 from the IV source
cone, and 18 from VI's parametric construction. The inherited IV modules
retain their source text and are credited in
\nolinkurl{AUTHORSHIP_NOTICE.md}. The note's own modules and gates carry
Apache-2.0 copyright headers; the complete license and attribution
notice are \nolinkurl{lean435/LICENSE} and \nolinkurl{lean435/NOTICE}. The
note's own modules and gates identify Juan Pablo Traverso Gianini as
author. The VI modules are synchronized with the author-credited source
version recorded in the manifest; existing third-party credits remain
unchanged.

The canonical root build checks \nolinkurl{ChordalStabilityNote}. It reuses
complete cached artifacts for unchanged dependencies, validates their
traces and recompiles changed sources. The source archive itself is
portable and contains no caches or local development junctions. This is
not a claim that all upstream modules or Mathlib were freshly rebuilt.

The full audit first verifies that every module in the source manifest
is present in the imported Lean environment; a missing module causes
rejection.

The selected root audit follows 7,125 declaration dependencies, requires
the structural and rounding inputs of IV and the parametric theorem of
VI, and excludes \nolinkurl{PaperIV.Erdos81Unconditional}. The full
local-project audit covers 7,393 declarations and visits 10,976
dependencies. Both permit only \nolinkurl{propext},
\nolinkurl{Classical.choice} and \nolinkurl{Quot.sound}; no additional
mathematical axiom or \nolinkurl{sorryAx} is accepted.
\nolinkurl{evidence/BUILD_CHECKS.json} records the final commands and
exits. \nolinkurl{FINAL_FORMAL_VALIDATION.json} binds them to the source
and public-entry-point checks.

\Needspace{8\baselineskip}
The immutable archive is
\nolinkurl{NOTE_II_LEAN435_FINAL_SOURCE_v1.0.zip}. Its SHA-256 is:

\begin{verbatim}
b24c5543d719779568ee9c099d453c170d4c69cbd4a31b293f4ee8375e287531
\end{verbatim}

\Needspace{8\baselineskip}
\nolinkurl{SOURCE_MANIFEST.json} records all source and configuration
hashes. Its SHA-256 is:

\begin{verbatim}
8bcdf425cee099e6e157dffa4356bb720e6fcd47ec138a2d4bac10b710587f75
\end{verbatim}

\nolinkurl{lean435/OWNERSHIP_AND_LICENSE.json} enumerates, with hashes,
all author-owned Lean sources, including the reused series cuts,
released by the author under Apache 2.0. The new full headers cover this
note's specific modules and audit entry points; inherited headers and
mathematical attributions are preserved. External dependencies retain
their own licenses.

\Needspace{12\baselineskip}
\subsection{Reproduction}\label{reproduction}

\Needspace{8\baselineskip}
Extract the archive and run:

\begin{verbatim}
cd lean435
lake exe cache get
lake build ChordalStabilityNote
lake env lean -j 1 -s 32768 Endpoint435Gate.lean
lake env lean -j 1 -s 32768 FreezeFullAudit.lean
\end{verbatim}

Preserve the supplied \nolinkurl{lake-manifest.json}; the cache command
retrieves artifacts for the pinned dependencies without requesting a new
dependency selection.

The root build executes the selected dependency and axiom audit. The
last command audits the complete imported local declaration set. The
vendored dependencies are included, so a checkout of the principal
papers is not required.

The endpoint retains \texttt{cp\ G\ \textless{}=\ Mval\ n} explicitly.
Without this hypothesis, natural subtraction could turn a negative
mathematical deficit into zero. In the manuscript this is the lower
bound on the deficit. The upper bound is
\texttt{2000\ *\ (Mval\ n\ -\ cp\ G)\ \textless{}=\ n}; integer division
does not weaken the window. Paper IV supplies the eventual upper bound
separately. Combining its threshold with this note's threshold gives the
corresponding written corollary for sufficiently large chordal graphs in
the stated window; the frozen endpoint keeps the explicit
nonnegative-deficit hypothesis.

\subsection{Manuscript correspondence and
scope}\label{manuscript-correspondence-and-scope}

The formal argument uses actual objects throughout. A completed far
packing is compared with the actual optimum. The near witness supplies
one regularized root with its physical packing and accounts. The
adapters identify those counts with VI's functions. The paid-root
theorem proves the entry window and applies the exact local
construction. Every field of \nolinkurl{RootStability} is a conclusion, and
the degree test identifies the same root at the end.

The Markdown sources \nolinkurl{CHORDAL_STABILITY_NOTE_v1.0_en.md} and
\nolinkurl{CHORDAL_STABILITY_NOTE_v1.0_es.md} preserve the theorem,
constants and displayed equations of the established derivation. The
current companion describes only the final validated source package. No
inverse transfer theorem, sharper threshold, new exceptional-vertex
extension or counting asymptotic is asserted.

\Needspace{18\baselineskip}\subsection{Assistance and
responsibility}\label{assistance-and-responsibility}

The principal note retains Paper IV's attribution of Claude, by
Anthropic, ChatGPT/Codex, by OpenAI, Aristotle, by Harmonic, Certo and
Jacobian for the inherited work. Aristotle's contribution to that work
included exploration, counterexample search, proof development,
formalization and review, and was not limited to translating finished
proofs. For this note, ChatGPT/Codex assisted with the argument, formal
interfaces, validation and presentation; Aristotle contributed a
separate formalization proposal. Exact-arithmetic regressions supplement
the source checks. Computations were reviewed within their stated scope;
a finite search does not replace a universal proof. Bibliographic
priority and external review are separate matters.

\Needspace{3\baselineskip}The author remains responsible for the arguments, citations, code, and
presentation. No AI system is listed as an author.

\end{document}
